\sectionkicker{Theme synthesis}
\subsection*{横断比較 — 科学Agentの価値は『完全自律』ではなくworkflow境界で読む}
\addcontentsline{toc}{subsection}{横断比較 — 科学Agentの価値は『完全自律』ではなくworkflow境界で読む}
6月のAI for Scienceでは、near-autonomous chemistry workflowとGPT-Rosalind更新を、単なるbenchmark向上ではなく、人間・model・実験操作・検証がどこで分担されるかというworkflow境界から読むことが重要である。
\medskip
\begin{center}
\small
\renewcommand{\arraystretch}{1.16}
\begin{tabularx}{\linewidth}{@{}p{0.20\linewidth}X@{}}
\toprule
観察軸 & 7月の一次資料・論文から読めること \\
\midrule
自律性の境界 & OAI-M1-03のchemistry workflowではhumanがproposal selection、lab operation、validationを担った。したがってnear-autonomousという表現をfully autonomousへ拡張しない。 \cite{src-5ff2f715f1e6} \\
実験結果の帰属 & tested substrate pairの11/14でyield improvementが報告されたが、これはOpenAI/Molecule.oneによる一次報告であり、一般的な化学実験性能へ拡張しない。 \cite{src-5ff2f715f1e6} \\
対象領域の広がり & GPT-RosalindのJune updateはmedicinal chemistry、genomics、lab-work能力をbenchmarksやvendor evaluationとともに提示した。chemistry workflowとは証拠形態が異なり、domain breadthの補助線として扱う。 \cite{src-ba562226725d} \\
workflowへの接続 & 2件を合わせると、6月の科学応用は『modelが科学問題へ回答する』だけでなく、proposal、実験計画、lab work、validationへどこまで接続できるかを問う段階へ進んでいる。ただし実世界の責任境界はhuman-in-the-loopとして残る。 \cite{src-5ff2f715f1e6,src-ba562226725d} \\
\bottomrule
\end{tabularx}
\renewcommand{\arraystretch}{1.0}
\end{center}
\normalsize
\begin{claimboundary}[この比較の境界]
この表は本号で既に検証・参照している一次資料と論文を横断比較の形へ再配置したものである。vendor / project / author が報告した性能値や優位性は独立再現済みの一般則へ変換せず、本文とSource Notesの帰属境界をそのまま維持する。
\end{claimboundary}
